\section{資料と未解決点}
\noindent 本号では、standard / specificationの公開状況、project・vendorの自己記述、研究結果を同じ強さの事実として扱わない。本文中のcitationから一次資料または論文へ戻れるようにしつつ、次の境界を残す。
\begin{itemize}
  \item IEEE 802.1DG、AUTOSAR、CAN XL等は公開されたstandard/project情報を確認したが、すべてのnormative clauseやparameterを規格本文まで再検証したわけではない。ISO 11898-1:2024の非公開範囲は推測で補わない。
  \item BMW、Mercedes-Benz、Volvo、Continental、DENSO/SOAFEE等のarchitecture・製品・platform説明は、各社・各projectの公式な自己記述として扱う。量産適用の表明と、独立したinteroperability / safety certification / performance validationは区別する。
  \item 研究論文の定量結果は、それぞれのhardware、testbed、task graph、simulation、laboratory条件に拘束される。異なるvehicle platformへ数値をそのまま移植しない。
  \item release dateやpublication dateの精度は公開資料の粒度を維持する。月単位までしか確認できない事象を人工的に日・時刻へ細分化しない。
  \item OSS release chronologyは実装の変化を追う手掛かりになるが、releaseの存在だけでproduction deploymentやconformanceを証明したとは扱わない。
\end{itemize}
\medskip
\noindent\textbf{References / Source Notes}\par
\smallskip
\noindent\small 共通の用途説明はここへ集約し、各entryではsource title、organization、date、URL、access dateを優先して表示する。詳細な検証状態やsourceごとのmachine-readable provenanceはsource bundle側に保持する。\normalsize
\printbibliography[heading=none]
